\section{Deep sequential computations with polynomial sample bounds}
\label{sec:critical-sequential-programs}

Some computations have long sequential implementations while admitting
sample bounds without a teacher-depth exponent. Decision trees and
binary-state branching programs have small Fourier variation. A different
argument handles permutation programs with solvable transition groups,
even when their Fourier variation is exponential. All results in this
section concern the original unrestricted critical prior and inherit the
sufficiently-large-width convention of
Theorem~\ref{thm:critical-excess-learning}.

\begin{proposition}[Decision trees and binary-state programs]
\label{prop:critical-binary-state-learning}
A Boolean function computed by a decision tree with $s$ leaves satisfies
$A_F(f)\le s$. A function computed by a deterministic layered branching
program of width at most two and length $T$ satisfies $A_F(f)\le T+1$.
The branching program may query different input coordinates in its two
states, and may query a coordinate repeatedly.

Consequently, there are universal $c,C>0$ such that a sufficient sample
bound for either class is
\[
 n\ge\frac C\varepsilon
       \left\{(d+L+K+2)^c+\log(1/\delta)\right\},
 \qquad L\ge C\log(d+L+K+2),
\]
where $K=s$ for the tree and $K=T+1$ for the branching program.
The guarantee is distribution-free on the cube, uses exact training fit,
and requires correct-label posterior probability greater than $2/3$
except on mass at most $\varepsilon$.
\end{proposition}

\begin{proof}
The indicator of a nonempty decision-tree leaf is a conjunction of input
literals, after removing repeated queries and contradictory branches.
Its Fourier $\ell^1$ norm is one. Expressing $f$ as the sum of its signed
leaf indicators proves the first assertion.

For the second, encode the two states by $\{-1,1\}$ and let $s_i(x)$
be the state after step $i$. Let $t_i^+(x)$ and $t_i^-(x)$ be the next
states selected from the two current states. Each is a unary Boolean
function of the coordinate queried at that state, hence a constant or a
signed input bit. Put
\[
 a_i=\frac{t_i^+-t_i^-}{2},\qquad
 b_i=\frac{t_i^++t_i^-}{2}.
\]
Then $s_i=a_is_{i-1}+b_i$, and the triangle inequality gives
$A_F(a_i),A_F(b_i)\le1$.
Fourier $\ell^1$ norm is submultiplicative: multiplication is convolution
on the group of subsets under symmetric difference. Therefore
\[
 A_F(s_i)\le A_F(s_{i-1})+1.
\]
The initial state is constant, so induction gives $A_F(s_T)\le T+1$.
The final output map is either constant or a signed copy of the state.
Neither case increases this bound. Apply
Corollary~\ref{cor:critical-fourier-learning}.
\end{proof}

Here width two counts discrete control states; it does not mean two
continuously valued neural activations. For a layered decision tree of
width $C$ and depth $h$, its number of leaves is at most $1+Ch$.
Thus the tree bound is jointly polynomial in its width and depth,
regardless of tree imbalance. Sharing in a general circuit can make a
decision-tree representation much larger.

\subsection{Solvable-group programs}

A length-$T$ program over a finite group $G$ chooses a group element
$g_i(x)$ at step $i$ according to one queried input bit, and outputs
\[
 f(x)=\phi(g_1(x)\cdots g_T(x)),\qquad
 \phi:G\longrightarrow\{-1,1\}.
\]
The queried coordinates may repeat. The group and the output table are
part of the specified implementation. The derived length of a solvable
group is the number of abelian quotients in its derived series.

\begin{lemma}[A uniform threshold compiler for solvable groups]
\label{lem:critical-solvable-group-compiler}
There is a universal $C$ such that the following holds. Suppose $G$ is
a finite solvable group of order $q\ge2$ and derived length $\ell\ge1$.
Every length-$T$ group program has a layered Boolean threshold
implementation with
\[
 h\le C\ell,\qquad
 C_T\le q^C(T+1)^C,\qquad
 B_T\le q^C(T+1)^C.
\]
Here $h$ is depth, $C_T$ is width including constant rails and layering
overhead, and $B_T$ bounds the Euclidean norms of gate coefficient
vectors. Gate scores have absolute value at least one on every
realized Boolean input. The constants are independent of the group.
\end{lemma}

\begin{proof}
We prove the stronger assertion that one circuit computes one-hot
encodings of every prefix product from one-hot encodings of the letters. A lookup involving at most three elements
of a group of order at most $q$ is a constant-depth Boolean circuit
of size $O(q^4)$: form the relevant conjunctions, then their disjunctions
for each output value. Its coefficients and fan-ins are polynomial in
$q$, with universal degrees.

For an abelian group, use a decomposition into cyclic groups of orders
$q_j\le q$, with at most $\log_2q$ factors. For one cyclic factor, let
$v_i\in\{0,\ldots,q_j-1\}$ be the input values. For every prefix length $i\le T$, form the
integer score
\[
 U_i=\sum_{r=1}^i\sum_{a=0}^{q_j-1}a\,\boldsymbol1_{\{v_r=a\}}.
\]
Compute $z_{i,s}=\boldsymbol1_{\{U_i\ge s\}}$ for
$1\le s\le(q_j-1)i$. For each residue $u$, put
$\psi_u(t)=2\boldsymbol1_{\{t\equiv u\pmod{q_j}\}}-1$.
On this integer interval the exact identity
\[
 \psi_u(U_i)=\psi_u(0)+
   \sum_{s=1}^{(q_j-1)i}[\psi_u(s)-\psi_u(s-1)]z_{i,s}
\]
shows that one further threshold layer computes the residue indicator.
For each prefix, combining the cyclic coordinates uses one conjunction
per group element: there are at most $q$ valid tuples, with at most
$\log_2q$ coordinates each. Computing all prefixes uses
$q^{O(1)}(T+1)^2$ gates. The size, fan-ins, and coefficients are
polynomial in $q,T$, with universal degrees, and the depth is
universally bounded. Half-integer threshold
offsets, followed by multiplication by two, give integer margins at
least one. Switching between $\{0,1\}$ and $\{-1,1\}$ conventions
preserves these polynomial bounds.

For the inductive step, take a nontrivial abelian normal subgroup
$A\triangleleft G$ such that $H=G/A$ has smaller derived length;
the last nontrivial term of the derived series suffices. Fix a section
$s:H\to G$ with $s(e)=e$, and write $g_i=a_i s(h_i)$ uniquely.
Set $H_i=h_1\cdots h_i$, with $H_0=e$. Define the finite tables
\[
 \alpha_h(a)=s(h)as(h)^{-1},\qquad
 c(h,h')=s(h)s(h')s(hh')^{-1}\in A.
\]
The group need not be a split extension. Compute all $H_i$ in parallel
using one copy of the stronger induction hypothesis, and then put
\[
 t_i=\alpha_{H_{i-1}}(a_i)c(H_{i-1},h_i)\in A.
\]
These are finite lookups. Multiplying the decompositions gives
\[
 g_1\cdots g_i=(t_1\cdots t_i)s(H_i),\qquad 1\le i\le T.
\]
Indeed a prefix $A_{i-1}s(H_{i-1})$ times $a_i s(h_i)$ equals
$A_{i-1}\alpha_{H_{i-1}}(a_i)c(H_{i-1},h_i)s(H_i)$.
All prefix products of the $t_i$ lie in the abelian group $A$ and are
computed simultaneously by the preceding construction. Final lookups
combine them with their corresponding $H_i$.

Each extension step increases depth by a universal constant and adds
$q^{O(1)}(T+1)^2$ gates to one quotient-prefix circuit. More explicitly,
for universal $c,C$, the raw size satisfies
\[
 F(G,T)\le F(G/A,T)+Cq^c(T+1)^2.
\]
The quotient order decreases by at least a factor two at every step,
so summing this geometric bound gives $F(G,T)\le Cq^c(T+1)^2$.
Every new gate has fan-in and coefficients polynomial in $q,T$ with
universal degrees. No step enlarges previously constructed coefficients.
The same uniform bound therefore holds for gate coefficient norms. Layering with identity gates and constant rails costs
only a polynomial in raw size and depth. Since $\ell\le\log_2q$, it is
absorbed in the stated width bound and does not increase coefficient
norms except by a universal polynomial factor.

Finally, choosing between two group letters requires only literals and
constants in their one-hot encoding. The output map $\phi$ is a finite
lookup. These additions preserve all bounds.
\end{proof}

\begin{theorem}[Learning solvable-group programs]
\label{thm:critical-solvable-group-learning}
There are universal $c,C>0$ such that the following holds. Let $f$ be
a length-$T$ program over a solvable group of order $q\ge2$ and derived
length $\ell\ge1$, and put $R=d+q+T+L+2$. If
\[
 L\ge C\ell\log R,
\]
a sufficient sample size for the original critical posterior is
\[
 n\ge\frac C\varepsilon
          \left\{R^{c\ell}+\log(1/\delta)\right\}.
\]
The prediction and width guarantees are those of
Theorem~\ref{thm:critical-excess-learning}. In particular, for bounded
derived length the bound is jointly polynomial in group order, program
length, and learner depth.
\end{theorem}

\begin{proof}
Apply Lemma~\ref{lem:critical-solvable-group-compiler}. The parameter in
Theorem~\ref{thm:controlled-circuit-rate} satisfies
\[
 T_c=\log(d+2)+h\log(\max\{d+1,C_T\}+B_T+h+L+2)
       \le C\ell\log R.
\]
The depth hypothesis ensures $L-h\ge L/2$ and $L-h\ge C_0T_c$ after
increasing the universal constant. Thus
\[
 J_L(f)\le V_*+C\frac{\ell^2\log^2R}{L}.
\]
The exponential factor in
Theorem~\ref{thm:critical-excess-learning} is at most $R^{C\ell}$.
Its remaining polynomial prefactor is absorbed in the same power,
since $\ell\le\log_2q$ and the displayed rate bound is polynomial in
$R$.
\end{proof}

\begin{corollary}[Permutation programs of width at most four]
\label{cor:critical-width-four-permutation-learning}
There are universal $c,C>0$ such that every length-$T$ permutation
branching program of width at most four has sufficient sample size
\[
 n\ge\frac C\varepsilon
       \left\{(d+T+L+2)^c+\log(1/\delta)\right\},
 \qquad L\ge C\log(d+T+L+2).
\]
Its original sequential depth $T$ is unrestricted.
\end{corollary}

\begin{proof}
A width-$w$ permutation program is a group program over $S_w$, with
its output determined by the action of the product on a fixed initial
state. For $w\le4$, the order and derived length are universally bounded.
The normal series
$\{e\}\triangleleft V_4\triangleleft A_4\triangleleft S_4$
has abelian successive quotients, so $S_4$ and its subgroups are solvable.
Apply the theorem. Width one gives a constant target, already covered
by Corollary~\ref{cor:critical-fourier-learning}.
\end{proof}

The permutation condition is material: the group proof does not cover
arbitrary noninvertible transformation programs of width three or four.
Nor does it prove a lower bound for nonsolvable groups. These are
restricted positive results for deep computations, and leave general
small-circuit learning unresolved.

\paragraph{Large Fourier variation can coexist with polynomial learning.}
For $b_i\in\{0,1\}$ and $x_i=(-1)^{b_i}$, consider the three-state
program
\[
 f(x)=2\boldsymbol1_{\{\sum_i b_i\equiv0\pmod3\}}-1.
\]
With $\alpha=(1+\sqrt3)/2>1$, it satisfies
\[
 A_F(f)\ge\tfrac13\alpha^d-\tfrac23.
\]
Indeed the roots-of-unity formula gives, for a nonempty Fourier set of
size $k$,
\[
 \widehat f(S)=\frac4{3\,2^d}3^{k/2}
                   \cos(d\pi/3-k\pi/2),
\]
with an additional $-1/3$ for the empty set. Retain the even $k$ terms,
use $|\cos(d\pi/3)|\ge1/2$, and apply the binomial theorem.
Nevertheless this is a program over the cyclic group of order three,
so Theorem~\ref{thm:critical-solvable-group-learning} gives polynomial
sample complexity at logarithmic learner depth. Large Fourier variation
alone is therefore insufficient for an adverse critical-prior learning
bound.
